\question
Однажды знаменитому Пэркюлю Уаро довелось расследовать удивительное дело о пропаже буквы $B$ и страшной смерти буквы $A$. Проведя весь день за допросами, Пэркюль смог определить главного подозреваемого, однако для окончательного вердикта пришлось записать в дневник следующее:
\\
\\
«Ни в коем случае не мог на мокрой трубе лежать нож, это взаимоисключающие события. Также обязательно должен был быть путь отступления, либо сообщник. Только тогда буква $B$ – преступник»
\\
\\
Помогите детективу, записав высказывание «Буква $B$ – преступник» в виде булевой функции, где $x_1$ – труба была мокрой, $x_2$ – на трубе лежал нож, $x_3$ – был путь отступления, $x_4$ – был сообщник. А также:
\begin{enumerate}
    \item Составьте СДНФ и СКНФ получившейся функции.
    \item Составьте полином Жегалкина получившейся функции любыми 2 способами.
\end{enumerate}

---------------

Автор -- Константин Васильев